\documentclass[11pt,a4paper]{article}
\usepackage[margin=23mm]{geometry}
\usepackage{amsmath,amssymb,amsthm}
\newtheorem{theorem}{Theorem}
\newcommand{\gen}[1]{\langle#1\rangle}
\title{A separator after the six-point head at odd degree}
\author{Human Compression project: L044 continuation}
\date{}
\begin{document}\maketitle
Let $G\le\operatorname{Alt}(\Omega)$ be transitive with connected
string C-marking $\rho_0,\ldots,\rho_{r-1}$, all maximal omissions
intransitive, $r\ge6$ and $|\Omega|=2r+1$. Its first three marks
are supported and transitive on $D=\{1,\ldots,6\}$. Assume the
actual interfaces
\[
 \rho_3=(6\ \beta)\delta,\qquad
 G_{\ge3}=\operatorname{Alt}(\Omega\setminus\{1,\ldots,5\}),
 \qquad \rho_4|_D=1,\qquad G_{\ge5}|_{D\cup\{\beta\}}=1.
\]
Here $\delta$ is an odd involution supported outside $D\cup\{\beta\}$.
For every $j\ge5$ assume there are two $j$-edges joining distinct
orbits of the full ambient omission $G_j$. They may lie in different
orbits of $H=G_{\ge5}$. Put $K=G_{>5}$.

\begin{theorem}
There is an index $s$ with
\[
                  5\le s<r-1,\qquad
                  |\operatorname{Fix}(G_{>s})|\ge2s+1.
\]
\end{theorem}
\begin{proof}
Use the complete strong simultaneous-witness theorem. It supplies
a colour $j\ge5$ with at most one ambient-eligible edge in each
actual $H$-orbit, and at least two distinct witness orbits. Every
witness on which $K$ is intransitive has first active label $5$,
exactly one edge eligible for its full local $j$-omission, and is
fixed pointwise by $\delta$.

If any witness is $K$-transitive, the complete transitive-witness
separator applies, with its actual whole-orbit transport. Its last
active $K$-label $x\ge6$ is nonterminal and gives the required
bound with $s=x$.

Otherwise select two distinct witnesses, and let their last active
labels be $x\le y$. Both contain the colour $j$, so $x\ge j$.
If $x=y$, the complete simultaneous-witness theorem gives
$x<r-1$ and at least $2x+1$ fixed points; take $s=x$.
If $x<y$ and $x\ge j+2$, apply the complete unequal-witness
separator with $s=x$. If $x=j+1<y$, use the complete one-step
separator with the same $s=x$.

The only remaining case is $x=j<y$. The complete crossing-endpoint
separator gives $s=j+1$ when $j\ge6$, and $s\in\{6,y\}$ when
$j=5$. In each case it explicitly proves $s<r-1$ and the stated
fixed-point bound. All its hypotheses have been supplied by the
simultaneous witness theorem. These cases exhaust every possible
pair of endpoint labels and every possible first-omission action.
\end{proof}

\paragraph{Attribution and scope.}
The actual head and orbit-deficit strategy originate in the proof of
Proposition 5.31 of Cameron--Fernandes--Leemans--Mixer,
\emph{Highest rank of a polytope for $A_n$}, arXiv:1605.09173v1.
This assembly uses the complete retained proofs of the five stated
separator results and their supporting orbit and parity arguments.
It adds only the exhaustive case routing; no new induction,
computation or formal statement is required.
The fixed points are those of the actual group in the displayed
theorem. Applying it to a proper prefix requires a separate proof
that the required fixed points survive any later attachment.
No complete L044 recovery or whole-original credit is claimed.
\end{document}
